% !TeX root = book.tex
% Glossary appendix draft for inclusion in the book, e.g.
% % !TeX root = book.tex
% Glossary appendix draft for inclusion in the book, e.g.
% % !TeX root = book.tex
% Glossary appendix draft for inclusion in the book, e.g.
% % !TeX root = book.tex
% Glossary appendix draft for inclusion in the book, e.g.
% \input{book/glossary}

\chapter{Glossary}\label{sec:glossary}

\noindent This appendix collects working definitions for core terms used across
the book and the current node implementation. The bias is conceptual first and
operational second: each entry states what the term means in the book and, when
useful, how the present implementation realizes it.

\paragraph{Actor.}
An actor is a long-lived concurrent process with a pid, a mailbox, and an
explicit message interface. In the ACTOR profile, actors are the basic unit of
concurrency. Toplevels, statechart processes, and node-resident services are all
specialized actor forms.

\paragraph{Agent.}
An agent is a persistent interactive process that participates in ongoing
interaction with its environment. In this book, a Prolog agent is defined by its
interface and behavior rather than by its implementation language: it talks
Prolog terms to other agents by receiving goals or messages and returning
answers, messages, or effects in the same term-based form.

Actors, toplevels, and nodes can all count as agents, but they play different
roles. Actors are loci of concurrent computation, toplevels expose a
conversational interface, and nodes host knowledge, expose web APIs, and enforce
policy.

\paragraph{Capability.}
A capability is authority to do or affect something.

In the book, the term is often used in the capability-security sense: possession
of an unforgeable or hard-to-guess handle, such as a pid or a reply channel,
confers some authority. In the current node implementation, the auth layer also
uses named policy capabilities such as \texttt{execute}, \texttt{admin}, and
\texttt{internal\_transport}, attached to a principal. These are different
mechanisms, but they answer the same question: what authority is being granted?

\paragraph{Containment.}
Containment is the structural security discipline that keeps client-created
computation bounded in lifetime and effect. A contained computation may still do
real work, but it should not outlive its session, escape its authority bounds,
or leave behind uncontrolled background activity.

In the current node, containment is reflected in the default rule that
client-spawned actors are linked to the client or session that created them.
Long-lived detached services are owner-managed exceptions, not the default.

\paragraph{Federation.}
Federation is the decision to standardize what systems share at the boundary
rather than to force them to become internally identical. In the Trinity, this
means shared web-facing interaction surfaces and execution profiles, while
leaving room for different Prolog systems to diverge internally.

\paragraph{Instant Portability.}
Instant portability is the limiting case of portability in which relocation
requires no adaptation step. A program can be downloaded from one node and run
locally, or shipped to another node, immediately because everything it depends
on is either bundled with it or guaranteed by the target profile.

\paragraph{Interaction Surface.}
An interaction surface is the externally visible contract through which systems
meet: routes, message formats, APIs, protocol rules, and profile promises. The
book uses this term to emphasize that interoperability depends on shared
boundaries, not on shared internals.

\paragraph{Interoperability.}
Interoperability is the ability of independently implemented systems to work
together correctly through shared protocols, representations, and profile
contracts. In this project, it is a stronger notion than mere syntactic
compatibility: two nodes interoperate when a client can rely on the same
observable behavior across them at the advertised profile boundary.

\paragraph{Inexpressibility.}
Inexpressibility is the structural security property that certain dangerous
operations cannot be expressed at all in client code. Rather than relying only
on runtime refusal, the exposed language fragment excludes operations such as
direct file I/O, raw operating-system access, or other ambient-authority
mechanisms that would be unsafe to hand to arbitrary clients.

In the book, this is the sandbox side of the security story. In the current
node, profile checks and sandbox policy are the implementation mechanisms that
approximate this boundary.

\paragraph{Invisibility.}
Invisibility is the structural security property that a client can interact only
with actors or services whose handles or names have been intentionally revealed.
It is the name-discipline side of security: what cannot be named cannot be
addressed.

In the book, opaque pids and owner-curated service publication are the key
mechanisms. In the current node, this also means that published services are not
the same thing as the ordinary client name registry.

\paragraph{Node.}
A node is a network-addressable host of execution, knowledge, and policy. A
node exposes Web Prolog APIs, advertises a profile, enforces authentication and
authorization policy, and may host both shared predicates and node-resident
actors.

\paragraph{Owner.}
An owner is the authority responsible for configuring a node's public surface
and long-lived resources. The owner role includes installing shared predicates,
publishing or withdrawing services, changing node settings and policy, and
taking responsibility for durable service lifecycle concerns such as supervision
and recovery.

In the current implementation, owner authority is realized through
administrative policy rather than through ordinary client capabilities. It is
best understood as a role in the security model, not merely as ``the client who
happens to connect first.''

\paragraph{Pid.}
A pid is an opaque process identifier for an actor. It is a routable handle, not
a descriptive name, and in security terms it often functions as a capability.
Clients should treat a pid as an implementation handle that may be passed around
or used for communication, but not as something that should be globally
enumerable.

\paragraph{Private Database.}
A private database is the actor-local body of clauses and facts associated with
one actor or toplevel. It forms part of that process's private workspace rather
than the node's shared knowledge, and it is not directly accessible to other
actors.

In the book, the private database is treated as a capability rather than as a
defining feature of agency: some agents use it heavily, while others remain
effectively stateless. In practical terms, it can be populated at spawn time by
\texttt{load\_*} options, it may shadow definitions from the shared database,
and updates to it affect only that actor.

\paragraph{Portability.}
Portability is the extent to which code, programs, or interaction patterns can
move between nodes or implementations without semantic breakage or manual
rewriting. In the book, portability is pursued by standardizing profiles and
their mandatory capabilities rather than by demanding total internal uniformity
across Prolog systems.

\paragraph{Principal.}
A principal is the identity on whose behalf a request is made. It is not the
same thing as a user, session, connection, or pid. A principal may be an
anonymous caller, a human user, a service account, or a trusted internal
transport identity.

In the current node implementation, principals are resolved by the auth layer
and mapped to policy capabilities. Authorization answers what a principal may
do; it does not by itself determine which pid or session is involved.

\paragraph{Profile.}
A profile is a deliberately scoped contract describing which language features,
APIs, and interaction patterns are available. Profiles let Web Prolog specify a
shared operational subset without turning Web Prolog into a separate fork of
Prolog.

In the current node implementation, the configured node profile is node-scoped,
not principal-scoped. It governs which routes and goal forms are in profile, and
public execution checks enforce that boundary.

\paragraph{Public Surface.}
A node's public surface is the set of predicates, services, routes, and policies
that it intentionally exposes to clients. It is smaller than the total internal
runtime and should be treated as an explicit design choice rather than as an
accident of what happens to be running.

\paragraph{Registered Name.}
A registered name is a symbolic name bound to an actor. Ordinary registered
names are convenience names for actor lookup and messaging; they are not
necessarily public services.

Under the current regime, ordinary client registration and published service
registration are separate. Client-side \texttt{register/2}, \texttt{whereis/2},
and \texttt{unregister/1} operate within a client-visible naming scope, while
published services use an owner-controlled service registry.

\paragraph{Reply Channel.}
A reply channel is a communication endpoint created for the purpose of receiving
one reply or a narrowly bounded class of replies. In capability terms, it should
carry reply authority only, not general authority to invoke the client.

The book treats reply channels as a way to avoid leaking a client's full pid to
an untrusted peer. The general principle is that communication authority should
be as narrow and short-lived as the protocol step that requires it.

\paragraph{Service.}
A service is a durable, shared conversational resource intentionally offered by
a node for repeated use by multiple clients. In the Prolog Web, services may be
predicate-based or actor-based, but node-resident actor services are especially
important because they provide shared state or shared coordination over time.

In the current implementation, a published actor service is not just any actor
with a convenient name. It is something the node owner explicitly publishes
under a stable service name.

\paragraph{Server.}
A server, in the sense used by the book's generic-behavior chapters, is a
stateful actor that maintains private state and serves clients through a stable
request-response message protocol. The generic server behavior factors out the
machinery of message routing, state threading, and reply handling so that the
application-specific logic can be supplied as a callback predicate.

This is an Erlang-style behavior: it exposes a protocol, not a predicate-level
logical interface, and it does not bind the caller's variables directly.

\paragraph{Service Discovery.}
Service discovery is the process by which clients learn which services a node
intentionally makes available. In the Prolog Web, discovery should be treated as
an owner-curated publication mechanism, not as process introspection.

The important distinction is that discovery resolves published service names
without revealing everything that happens to be running. In the current book
example, this is realized declaratively through relations such as
\texttt{service/2} queryable via \texttt{rpc/2-3}.

\paragraph{Service Publication.}
Service publication is the act by which a node owner places a name into the
node's public service namespace. Publication is stronger than ordinary naming:
it makes a service part of the node's advertised public surface.

In the current regime, this is distinct from ordinary \texttt{register/2}.
Client registration is for client-managed actors, while published services are
managed through the owner-controlled service registry.

\paragraph{Session.}
A session is a bounded continuing interaction context in which state can persist
across multiple protocol steps. The book contrasts session-oriented interaction
with purely stateless request-response exchange.

In the current node, HTTP ISOTOPE sessions and WebSocket ACTOR sessions are both
session forms, but they have different timing and failure boundaries.

\paragraph{Shared Database.}
The shared database is the node-owned body of predicates and facts made
available as shared knowledge. It is loaded once at node startup into a shared
runtime module and then imported by actor modules or used by node-level query
execution.

Clients may query against shared knowledge and combine it with private code, but
the shared database is not the same thing as a client's private mutable
workspace.

\paragraph{Statechart Actor.}
A statechart actor is a Web Prolog actor whose control logic is given by a
statechart document. It behaves like an ordinary actor with a mailbox, but
incoming messages are interpreted as events that drive a structured state
machine rather than an ad hoc \texttt{receive/1-2} loop.

The distinctive semantic idea is run-to-completion: processing one external
event may trigger a finite cascade of internal work before the next external
message is consumed. This makes control explicit and inspectable while keeping
the process inside the ordinary actor model.

\paragraph{Supervisor.}
A supervisor is an actor whose main responsibility is to start, monitor, and
manage child actors, especially in the face of failure. In the book's Erlang-style
account, a supervisor embodies the ``let it crash'' discipline: workers are
allowed to fail, and the supervisor applies the configured recovery policy.

Supervisors matter most for durable node-resident services, where restart and
cleanup are part of the node owner's operational responsibility rather than
something each worker should reimplement ad hoc.

\paragraph{Toplevel Actor.}
A toplevel actor is an actor whose behavior is governed by an explicit
conversational protocol for submitting goals, receiving answers incrementally,
handling output, and managing abort or stop. It turns the ordinary Prolog
toplevel into a first-class actor process.

In implementation terms, toplevel actors underpin \texttt{rpc/2-3}, the
semi-stateful HTTP API, and much of the browser-facing shell machinery.

\paragraph{PTCP.}
PTCP is the Prolog Toplevel Communication Protocol. It is the explicit protocol
governing conversation between a client and a toplevel actor: how goals are
submitted, how answers are streamed incrementally, how output and prompts may
occur, and how abort, stop, and exit are handled.

Its importance in the book is that it turns informal REPL behavior into a
first-class behavioral contract. In implementation terms, PTCP underlies the
observable behavior of toplevel actors across the stateless, semi-stateful, and
stateful APIs.


\chapter{Glossary}\label{sec:glossary}

\noindent This appendix collects working definitions for core terms used across
the book and the current node implementation. The bias is conceptual first and
operational second: each entry states what the term means in the book and, when
useful, how the present implementation realizes it.

\paragraph{Actor.}
An actor is a long-lived concurrent process with a pid, a mailbox, and an
explicit message interface. In the ACTOR profile, actors are the basic unit of
concurrency. Toplevels, statechart processes, and node-resident services are all
specialized actor forms.

\paragraph{Agent.}
An agent is a persistent interactive process that participates in ongoing
interaction with its environment. In this book, a Prolog agent is defined by its
interface and behavior rather than by its implementation language: it talks
Prolog terms to other agents by receiving goals or messages and returning
answers, messages, or effects in the same term-based form.

Actors, toplevels, and nodes can all count as agents, but they play different
roles. Actors are loci of concurrent computation, toplevels expose a
conversational interface, and nodes host knowledge, expose web APIs, and enforce
policy.

\paragraph{Capability.}
A capability is authority to do or affect something.

In the book, the term is often used in the capability-security sense: possession
of an unforgeable or hard-to-guess handle, such as a pid or a reply channel,
confers some authority. In the current node implementation, the auth layer also
uses named policy capabilities such as \texttt{execute}, \texttt{admin}, and
\texttt{internal\_transport}, attached to a principal. These are different
mechanisms, but they answer the same question: what authority is being granted?

\paragraph{Containment.}
Containment is the structural security discipline that keeps client-created
computation bounded in lifetime and effect. A contained computation may still do
real work, but it should not outlive its session, escape its authority bounds,
or leave behind uncontrolled background activity.

In the current node, containment is reflected in the default rule that
client-spawned actors are linked to the client or session that created them.
Long-lived detached services are owner-managed exceptions, not the default.

\paragraph{Federation.}
Federation is the decision to standardize what systems share at the boundary
rather than to force them to become internally identical. In the Trinity, this
means shared web-facing interaction surfaces and execution profiles, while
leaving room for different Prolog systems to diverge internally.

\paragraph{Instant Portability.}
Instant portability is the limiting case of portability in which relocation
requires no adaptation step. A program can be downloaded from one node and run
locally, or shipped to another node, immediately because everything it depends
on is either bundled with it or guaranteed by the target profile.

\paragraph{Interaction Surface.}
An interaction surface is the externally visible contract through which systems
meet: routes, message formats, APIs, protocol rules, and profile promises. The
book uses this term to emphasize that interoperability depends on shared
boundaries, not on shared internals.

\paragraph{Interoperability.}
Interoperability is the ability of independently implemented systems to work
together correctly through shared protocols, representations, and profile
contracts. In this project, it is a stronger notion than mere syntactic
compatibility: two nodes interoperate when a client can rely on the same
observable behavior across them at the advertised profile boundary.

\paragraph{Inexpressibility.}
Inexpressibility is the structural security property that certain dangerous
operations cannot be expressed at all in client code. Rather than relying only
on runtime refusal, the exposed language fragment excludes operations such as
direct file I/O, raw operating-system access, or other ambient-authority
mechanisms that would be unsafe to hand to arbitrary clients.

In the book, this is the sandbox side of the security story. In the current
node, profile checks and sandbox policy are the implementation mechanisms that
approximate this boundary.

\paragraph{Invisibility.}
Invisibility is the structural security property that a client can interact only
with actors or services whose handles or names have been intentionally revealed.
It is the name-discipline side of security: what cannot be named cannot be
addressed.

In the book, opaque pids and owner-curated service publication are the key
mechanisms. In the current node, this also means that published services are not
the same thing as the ordinary client name registry.

\paragraph{Node.}
A node is a network-addressable host of execution, knowledge, and policy. A
node exposes Web Prolog APIs, advertises a profile, enforces authentication and
authorization policy, and may host both shared predicates and node-resident
actors.

\paragraph{Owner.}
An owner is the authority responsible for configuring a node's public surface
and long-lived resources. The owner role includes installing shared predicates,
publishing or withdrawing services, changing node settings and policy, and
taking responsibility for durable service lifecycle concerns such as supervision
and recovery.

In the current implementation, owner authority is realized through
administrative policy rather than through ordinary client capabilities. It is
best understood as a role in the security model, not merely as ``the client who
happens to connect first.''

\paragraph{Pid.}
A pid is an opaque process identifier for an actor. It is a routable handle, not
a descriptive name, and in security terms it often functions as a capability.
Clients should treat a pid as an implementation handle that may be passed around
or used for communication, but not as something that should be globally
enumerable.

\paragraph{Private Database.}
A private database is the actor-local body of clauses and facts associated with
one actor or toplevel. It forms part of that process's private workspace rather
than the node's shared knowledge, and it is not directly accessible to other
actors.

In the book, the private database is treated as a capability rather than as a
defining feature of agency: some agents use it heavily, while others remain
effectively stateless. In practical terms, it can be populated at spawn time by
\texttt{load\_*} options, it may shadow definitions from the shared database,
and updates to it affect only that actor.

\paragraph{Portability.}
Portability is the extent to which code, programs, or interaction patterns can
move between nodes or implementations without semantic breakage or manual
rewriting. In the book, portability is pursued by standardizing profiles and
their mandatory capabilities rather than by demanding total internal uniformity
across Prolog systems.

\paragraph{Principal.}
A principal is the identity on whose behalf a request is made. It is not the
same thing as a user, session, connection, or pid. A principal may be an
anonymous caller, a human user, a service account, or a trusted internal
transport identity.

In the current node implementation, principals are resolved by the auth layer
and mapped to policy capabilities. Authorization answers what a principal may
do; it does not by itself determine which pid or session is involved.

\paragraph{Profile.}
A profile is a deliberately scoped contract describing which language features,
APIs, and interaction patterns are available. Profiles let Web Prolog specify a
shared operational subset without turning Web Prolog into a separate fork of
Prolog.

In the current node implementation, the configured node profile is node-scoped,
not principal-scoped. It governs which routes and goal forms are in profile, and
public execution checks enforce that boundary.

\paragraph{Public Surface.}
A node's public surface is the set of predicates, services, routes, and policies
that it intentionally exposes to clients. It is smaller than the total internal
runtime and should be treated as an explicit design choice rather than as an
accident of what happens to be running.

\paragraph{Registered Name.}
A registered name is a symbolic name bound to an actor. Ordinary registered
names are convenience names for actor lookup and messaging; they are not
necessarily public services.

Under the current regime, ordinary client registration and published service
registration are separate. Client-side \texttt{register/2}, \texttt{whereis/2},
and \texttt{unregister/1} operate within a client-visible naming scope, while
published services use an owner-controlled service registry.

\paragraph{Reply Channel.}
A reply channel is a communication endpoint created for the purpose of receiving
one reply or a narrowly bounded class of replies. In capability terms, it should
carry reply authority only, not general authority to invoke the client.

The book treats reply channels as a way to avoid leaking a client's full pid to
an untrusted peer. The general principle is that communication authority should
be as narrow and short-lived as the protocol step that requires it.

\paragraph{Service.}
A service is a durable, shared conversational resource intentionally offered by
a node for repeated use by multiple clients. In the Prolog Web, services may be
predicate-based or actor-based, but node-resident actor services are especially
important because they provide shared state or shared coordination over time.

In the current implementation, a published actor service is not just any actor
with a convenient name. It is something the node owner explicitly publishes
under a stable service name.

\paragraph{Server.}
A server, in the sense used by the book's generic-behavior chapters, is a
stateful actor that maintains private state and serves clients through a stable
request-response message protocol. The generic server behavior factors out the
machinery of message routing, state threading, and reply handling so that the
application-specific logic can be supplied as a callback predicate.

This is an Erlang-style behavior: it exposes a protocol, not a predicate-level
logical interface, and it does not bind the caller's variables directly.

\paragraph{Service Discovery.}
Service discovery is the process by which clients learn which services a node
intentionally makes available. In the Prolog Web, discovery should be treated as
an owner-curated publication mechanism, not as process introspection.

The important distinction is that discovery resolves published service names
without revealing everything that happens to be running. In the current book
example, this is realized declaratively through relations such as
\texttt{service/2} queryable via \texttt{rpc/2-3}.

\paragraph{Service Publication.}
Service publication is the act by which a node owner places a name into the
node's public service namespace. Publication is stronger than ordinary naming:
it makes a service part of the node's advertised public surface.

In the current regime, this is distinct from ordinary \texttt{register/2}.
Client registration is for client-managed actors, while published services are
managed through the owner-controlled service registry.

\paragraph{Session.}
A session is a bounded continuing interaction context in which state can persist
across multiple protocol steps. The book contrasts session-oriented interaction
with purely stateless request-response exchange.

In the current node, HTTP ISOTOPE sessions and WebSocket ACTOR sessions are both
session forms, but they have different timing and failure boundaries.

\paragraph{Shared Database.}
The shared database is the node-owned body of predicates and facts made
available as shared knowledge. It is loaded once at node startup into a shared
runtime module and then imported by actor modules or used by node-level query
execution.

Clients may query against shared knowledge and combine it with private code, but
the shared database is not the same thing as a client's private mutable
workspace.

\paragraph{Statechart Actor.}
A statechart actor is a Web Prolog actor whose control logic is given by a
statechart document. It behaves like an ordinary actor with a mailbox, but
incoming messages are interpreted as events that drive a structured state
machine rather than an ad hoc \texttt{receive/1-2} loop.

The distinctive semantic idea is run-to-completion: processing one external
event may trigger a finite cascade of internal work before the next external
message is consumed. This makes control explicit and inspectable while keeping
the process inside the ordinary actor model.

\paragraph{Supervisor.}
A supervisor is an actor whose main responsibility is to start, monitor, and
manage child actors, especially in the face of failure. In the book's Erlang-style
account, a supervisor embodies the ``let it crash'' discipline: workers are
allowed to fail, and the supervisor applies the configured recovery policy.

Supervisors matter most for durable node-resident services, where restart and
cleanup are part of the node owner's operational responsibility rather than
something each worker should reimplement ad hoc.

\paragraph{Toplevel Actor.}
A toplevel actor is an actor whose behavior is governed by an explicit
conversational protocol for submitting goals, receiving answers incrementally,
handling output, and managing abort or stop. It turns the ordinary Prolog
toplevel into a first-class actor process.

In implementation terms, toplevel actors underpin \texttt{rpc/2-3}, the
semi-stateful HTTP API, and much of the browser-facing shell machinery.

\paragraph{PTCP.}
PTCP is the Prolog Toplevel Communication Protocol. It is the explicit protocol
governing conversation between a client and a toplevel actor: how goals are
submitted, how answers are streamed incrementally, how output and prompts may
occur, and how abort, stop, and exit are handled.

Its importance in the book is that it turns informal REPL behavior into a
first-class behavioral contract. In implementation terms, PTCP underlies the
observable behavior of toplevel actors across the stateless, semi-stateful, and
stateful APIs.


\chapter{Glossary}\label{sec:glossary}

\noindent This appendix collects working definitions for core terms used across
the book and the current node implementation. The bias is conceptual first and
operational second: each entry states what the term means in the book and, when
useful, how the present implementation realizes it.

\paragraph{Actor.}
An actor is a long-lived concurrent process with a pid, a mailbox, and an
explicit message interface. In the ACTOR profile, actors are the basic unit of
concurrency. Toplevels, statechart processes, and node-resident services are all
specialized actor forms.

\paragraph{Agent.}
An agent is a persistent interactive process that participates in ongoing
interaction with its environment. In this book, a Prolog agent is defined by its
interface and behavior rather than by its implementation language: it talks
Prolog terms to other agents by receiving goals or messages and returning
answers, messages, or effects in the same term-based form.

Actors, toplevels, and nodes can all count as agents, but they play different
roles. Actors are loci of concurrent computation, toplevels expose a
conversational interface, and nodes host knowledge, expose web APIs, and enforce
policy.

\paragraph{Capability.}
A capability is authority to do or affect something.

In the book, the term is often used in the capability-security sense: possession
of an unforgeable or hard-to-guess handle, such as a pid or a reply channel,
confers some authority. In the current node implementation, the auth layer also
uses named policy capabilities such as \texttt{execute}, \texttt{admin}, and
\texttt{internal\_transport}, attached to a principal. These are different
mechanisms, but they answer the same question: what authority is being granted?

\paragraph{Containment.}
Containment is the structural security discipline that keeps client-created
computation bounded in lifetime and effect. A contained computation may still do
real work, but it should not outlive its session, escape its authority bounds,
or leave behind uncontrolled background activity.

In the current node, containment is reflected in the default rule that
client-spawned actors are linked to the client or session that created them.
Long-lived detached services are owner-managed exceptions, not the default.

\paragraph{Federation.}
Federation is the decision to standardize what systems share at the boundary
rather than to force them to become internally identical. In the Trinity, this
means shared web-facing interaction surfaces and execution profiles, while
leaving room for different Prolog systems to diverge internally.

\paragraph{Instant Portability.}
Instant portability is the limiting case of portability in which relocation
requires no adaptation step. A program can be downloaded from one node and run
locally, or shipped to another node, immediately because everything it depends
on is either bundled with it or guaranteed by the target profile.

\paragraph{Interaction Surface.}
An interaction surface is the externally visible contract through which systems
meet: routes, message formats, APIs, protocol rules, and profile promises. The
book uses this term to emphasize that interoperability depends on shared
boundaries, not on shared internals.

\paragraph{Interoperability.}
Interoperability is the ability of independently implemented systems to work
together correctly through shared protocols, representations, and profile
contracts. In this project, it is a stronger notion than mere syntactic
compatibility: two nodes interoperate when a client can rely on the same
observable behavior across them at the advertised profile boundary.

\paragraph{Inexpressibility.}
Inexpressibility is the structural security property that certain dangerous
operations cannot be expressed at all in client code. Rather than relying only
on runtime refusal, the exposed language fragment excludes operations such as
direct file I/O, raw operating-system access, or other ambient-authority
mechanisms that would be unsafe to hand to arbitrary clients.

In the book, this is the sandbox side of the security story. In the current
node, profile checks and sandbox policy are the implementation mechanisms that
approximate this boundary.

\paragraph{Invisibility.}
Invisibility is the structural security property that a client can interact only
with actors or services whose handles or names have been intentionally revealed.
It is the name-discipline side of security: what cannot be named cannot be
addressed.

In the book, opaque pids and owner-curated service publication are the key
mechanisms. In the current node, this also means that published services are not
the same thing as the ordinary client name registry.

\paragraph{Node.}
A node is a network-addressable host of execution, knowledge, and policy. A
node exposes Web Prolog APIs, advertises a profile, enforces authentication and
authorization policy, and may host both shared predicates and node-resident
actors.

\paragraph{Owner.}
An owner is the authority responsible for configuring a node's public surface
and long-lived resources. The owner role includes installing shared predicates,
publishing or withdrawing services, changing node settings and policy, and
taking responsibility for durable service lifecycle concerns such as supervision
and recovery.

In the current implementation, owner authority is realized through
administrative policy rather than through ordinary client capabilities. It is
best understood as a role in the security model, not merely as ``the client who
happens to connect first.''

\paragraph{Pid.}
A pid is an opaque process identifier for an actor. It is a routable handle, not
a descriptive name, and in security terms it often functions as a capability.
Clients should treat a pid as an implementation handle that may be passed around
or used for communication, but not as something that should be globally
enumerable.

\paragraph{Private Database.}
A private database is the actor-local body of clauses and facts associated with
one actor or toplevel. It forms part of that process's private workspace rather
than the node's shared knowledge, and it is not directly accessible to other
actors.

In the book, the private database is treated as a capability rather than as a
defining feature of agency: some agents use it heavily, while others remain
effectively stateless. In practical terms, it can be populated at spawn time by
\texttt{load\_*} options, it may shadow definitions from the shared database,
and updates to it affect only that actor.

\paragraph{Portability.}
Portability is the extent to which code, programs, or interaction patterns can
move between nodes or implementations without semantic breakage or manual
rewriting. In the book, portability is pursued by standardizing profiles and
their mandatory capabilities rather than by demanding total internal uniformity
across Prolog systems.

\paragraph{Principal.}
A principal is the identity on whose behalf a request is made. It is not the
same thing as a user, session, connection, or pid. A principal may be an
anonymous caller, a human user, a service account, or a trusted internal
transport identity.

In the current node implementation, principals are resolved by the auth layer
and mapped to policy capabilities. Authorization answers what a principal may
do; it does not by itself determine which pid or session is involved.

\paragraph{Profile.}
A profile is a deliberately scoped contract describing which language features,
APIs, and interaction patterns are available. Profiles let Web Prolog specify a
shared operational subset without turning Web Prolog into a separate fork of
Prolog.

In the current node implementation, the configured node profile is node-scoped,
not principal-scoped. It governs which routes and goal forms are in profile, and
public execution checks enforce that boundary.

\paragraph{Public Surface.}
A node's public surface is the set of predicates, services, routes, and policies
that it intentionally exposes to clients. It is smaller than the total internal
runtime and should be treated as an explicit design choice rather than as an
accident of what happens to be running.

\paragraph{Registered Name.}
A registered name is a symbolic name bound to an actor. Ordinary registered
names are convenience names for actor lookup and messaging; they are not
necessarily public services.

Under the current regime, ordinary client registration and published service
registration are separate. Client-side \texttt{register/2}, \texttt{whereis/2},
and \texttt{unregister/1} operate within a client-visible naming scope, while
published services use an owner-controlled service registry.

\paragraph{Reply Channel.}
A reply channel is a communication endpoint created for the purpose of receiving
one reply or a narrowly bounded class of replies. In capability terms, it should
carry reply authority only, not general authority to invoke the client.

The book treats reply channels as a way to avoid leaking a client's full pid to
an untrusted peer. The general principle is that communication authority should
be as narrow and short-lived as the protocol step that requires it.

\paragraph{Service.}
A service is a durable, shared conversational resource intentionally offered by
a node for repeated use by multiple clients. In the Prolog Web, services may be
predicate-based or actor-based, but node-resident actor services are especially
important because they provide shared state or shared coordination over time.

In the current implementation, a published actor service is not just any actor
with a convenient name. It is something the node owner explicitly publishes
under a stable service name.

\paragraph{Server.}
A server, in the sense used by the book's generic-behavior chapters, is a
stateful actor that maintains private state and serves clients through a stable
request-response message protocol. The generic server behavior factors out the
machinery of message routing, state threading, and reply handling so that the
application-specific logic can be supplied as a callback predicate.

This is an Erlang-style behavior: it exposes a protocol, not a predicate-level
logical interface, and it does not bind the caller's variables directly.

\paragraph{Service Discovery.}
Service discovery is the process by which clients learn which services a node
intentionally makes available. In the Prolog Web, discovery should be treated as
an owner-curated publication mechanism, not as process introspection.

The important distinction is that discovery resolves published service names
without revealing everything that happens to be running. In the current book
example, this is realized declaratively through relations such as
\texttt{service/2} queryable via \texttt{rpc/2-3}.

\paragraph{Service Publication.}
Service publication is the act by which a node owner places a name into the
node's public service namespace. Publication is stronger than ordinary naming:
it makes a service part of the node's advertised public surface.

In the current regime, this is distinct from ordinary \texttt{register/2}.
Client registration is for client-managed actors, while published services are
managed through the owner-controlled service registry.

\paragraph{Session.}
A session is a bounded continuing interaction context in which state can persist
across multiple protocol steps. The book contrasts session-oriented interaction
with purely stateless request-response exchange.

In the current node, HTTP ISOTOPE sessions and WebSocket ACTOR sessions are both
session forms, but they have different timing and failure boundaries.

\paragraph{Shared Database.}
The shared database is the node-owned body of predicates and facts made
available as shared knowledge. It is loaded once at node startup into a shared
runtime module and then imported by actor modules or used by node-level query
execution.

Clients may query against shared knowledge and combine it with private code, but
the shared database is not the same thing as a client's private mutable
workspace.

\paragraph{Statechart Actor.}
A statechart actor is a Web Prolog actor whose control logic is given by a
statechart document. It behaves like an ordinary actor with a mailbox, but
incoming messages are interpreted as events that drive a structured state
machine rather than an ad hoc \texttt{receive/1-2} loop.

The distinctive semantic idea is run-to-completion: processing one external
event may trigger a finite cascade of internal work before the next external
message is consumed. This makes control explicit and inspectable while keeping
the process inside the ordinary actor model.

\paragraph{Supervisor.}
A supervisor is an actor whose main responsibility is to start, monitor, and
manage child actors, especially in the face of failure. In the book's Erlang-style
account, a supervisor embodies the ``let it crash'' discipline: workers are
allowed to fail, and the supervisor applies the configured recovery policy.

Supervisors matter most for durable node-resident services, where restart and
cleanup are part of the node owner's operational responsibility rather than
something each worker should reimplement ad hoc.

\paragraph{Toplevel Actor.}
A toplevel actor is an actor whose behavior is governed by an explicit
conversational protocol for submitting goals, receiving answers incrementally,
handling output, and managing abort or stop. It turns the ordinary Prolog
toplevel into a first-class actor process.

In implementation terms, toplevel actors underpin \texttt{rpc/2-3}, the
semi-stateful HTTP API, and much of the browser-facing shell machinery.

\paragraph{PTCP.}
PTCP is the Prolog Toplevel Communication Protocol. It is the explicit protocol
governing conversation between a client and a toplevel actor: how goals are
submitted, how answers are streamed incrementally, how output and prompts may
occur, and how abort, stop, and exit are handled.

Its importance in the book is that it turns informal REPL behavior into a
first-class behavioral contract. In implementation terms, PTCP underlies the
observable behavior of toplevel actors across the stateless, semi-stateful, and
stateful APIs.


\chapter{Glossary}\label{sec:glossary}

\noindent This appendix collects working definitions for core terms used across
the book and the current node implementation. The bias is conceptual first and
operational second: each entry states what the term means in the book and, when
useful, how the present implementation realizes it.

\paragraph{Actor.}
An actor is a long-lived concurrent process with a pid, a mailbox, and an
explicit message interface. In the ACTOR profile, actors are the basic unit of
concurrency. Toplevels, statechart processes, and node-resident services are all
specialized actor forms.

\paragraph{Agent.}
An agent is a persistent interactive process that participates in ongoing
interaction with its environment. In this book, a Prolog agent is defined by its
interface and behavior rather than by its implementation language: it talks
Prolog terms to other agents by receiving goals or messages and returning
answers, messages, or effects in the same term-based form.

Actors, toplevels, and nodes can all count as agents, but they play different
roles. Actors are loci of concurrent computation, toplevels expose a
conversational interface, and nodes host knowledge, expose web APIs, and enforce
policy.

\paragraph{Capability.}
A capability is authority to do or affect something.

In the book, the term is often used in the capability-security sense: possession
of an unforgeable or hard-to-guess handle, such as a pid or a reply channel,
confers some authority. In the current node implementation, the auth layer also
uses named policy capabilities such as \texttt{execute}, \texttt{admin}, and
\texttt{internal\_transport}, attached to a principal. These are different
mechanisms, but they answer the same question: what authority is being granted?

\paragraph{Containment.}
Containment is the structural security discipline that keeps client-created
computation bounded in lifetime and effect. A contained computation may still do
real work, but it should not outlive its session, escape its authority bounds,
or leave behind uncontrolled background activity.

In the current node, containment is reflected in the default rule that
client-spawned actors are linked to the client or session that created them.
Long-lived detached services are owner-managed exceptions, not the default.

\paragraph{Federation.}
Federation is the decision to standardize what systems share at the boundary
rather than to force them to become internally identical. In the Trinity, this
means shared web-facing interaction surfaces and execution profiles, while
leaving room for different Prolog systems to diverge internally.

\paragraph{Instant Portability.}
Instant portability is the limiting case of portability in which relocation
requires no adaptation step. A program can be downloaded from one node and run
locally, or shipped to another node, immediately because everything it depends
on is either bundled with it or guaranteed by the target profile.

\paragraph{Interaction Surface.}
An interaction surface is the externally visible contract through which systems
meet: routes, message formats, APIs, protocol rules, and profile promises. The
book uses this term to emphasize that interoperability depends on shared
boundaries, not on shared internals.

\paragraph{Interoperability.}
Interoperability is the ability of independently implemented systems to work
together correctly through shared protocols, representations, and profile
contracts. In this project, it is a stronger notion than mere syntactic
compatibility: two nodes interoperate when a client can rely on the same
observable behavior across them at the advertised profile boundary.

\paragraph{Inexpressibility.}
Inexpressibility is the structural security property that certain dangerous
operations cannot be expressed at all in client code. Rather than relying only
on runtime refusal, the exposed language fragment excludes operations such as
direct file I/O, raw operating-system access, or other ambient-authority
mechanisms that would be unsafe to hand to arbitrary clients.

In the book, this is the sandbox side of the security story. In the current
node, profile checks and sandbox policy are the implementation mechanisms that
approximate this boundary.

\paragraph{Invisibility.}
Invisibility is the structural security property that a client can interact only
with actors or services whose handles or names have been intentionally revealed.
It is the name-discipline side of security: what cannot be named cannot be
addressed.

In the book, opaque pids and owner-curated service publication are the key
mechanisms. In the current node, this also means that published services are not
the same thing as the ordinary client name registry.

\paragraph{Node.}
A node is a network-addressable host of execution, knowledge, and policy. A
node exposes Web Prolog APIs, advertises a profile, enforces authentication and
authorization policy, and may host both shared predicates and node-resident
actors.

\paragraph{Owner.}
An owner is the authority responsible for configuring a node's public surface
and long-lived resources. The owner role includes installing shared predicates,
publishing or withdrawing services, changing node settings and policy, and
taking responsibility for durable service lifecycle concerns such as supervision
and recovery.

In the current implementation, owner authority is realized through
administrative policy rather than through ordinary client capabilities. It is
best understood as a role in the security model, not merely as ``the client who
happens to connect first.''

\paragraph{Pid.}
A pid is an opaque process identifier for an actor. It is a routable handle, not
a descriptive name, and in security terms it often functions as a capability.
Clients should treat a pid as an implementation handle that may be passed around
or used for communication, but not as something that should be globally
enumerable.

\paragraph{Private Database.}
A private database is the actor-local body of clauses and facts associated with
one actor or toplevel. It forms part of that process's private workspace rather
than the node's shared knowledge, and it is not directly accessible to other
actors.

In the book, the private database is treated as a capability rather than as a
defining feature of agency: some agents use it heavily, while others remain
effectively stateless. In practical terms, it can be populated at spawn time by
\texttt{load\_*} options, it may shadow definitions from the shared database,
and updates to it affect only that actor.

\paragraph{Portability.}
Portability is the extent to which code, programs, or interaction patterns can
move between nodes or implementations without semantic breakage or manual
rewriting. In the book, portability is pursued by standardizing profiles and
their mandatory capabilities rather than by demanding total internal uniformity
across Prolog systems.

\paragraph{Principal.}
A principal is the identity on whose behalf a request is made. It is not the
same thing as a user, session, connection, or pid. A principal may be an
anonymous caller, a human user, a service account, or a trusted internal
transport identity.

In the current node implementation, principals are resolved by the auth layer
and mapped to policy capabilities. Authorization answers what a principal may
do; it does not by itself determine which pid or session is involved.

\paragraph{Profile.}
A profile is a deliberately scoped contract describing which language features,
APIs, and interaction patterns are available. Profiles let Web Prolog specify a
shared operational subset without turning Web Prolog into a separate fork of
Prolog.

In the current node implementation, the configured node profile is node-scoped,
not principal-scoped. It governs which routes and goal forms are in profile, and
public execution checks enforce that boundary.

\paragraph{Public Surface.}
A node's public surface is the set of predicates, services, routes, and policies
that it intentionally exposes to clients. It is smaller than the total internal
runtime and should be treated as an explicit design choice rather than as an
accident of what happens to be running.

\paragraph{Registered Name.}
A registered name is a symbolic name bound to an actor. Ordinary registered
names are convenience names for actor lookup and messaging; they are not
necessarily public services.

Under the current regime, ordinary client registration and published service
registration are separate. Client-side \texttt{register/2}, \texttt{whereis/2},
and \texttt{unregister/1} operate within a client-visible naming scope, while
published services use an owner-controlled service registry.

\paragraph{Reply Channel.}
A reply channel is a communication endpoint created for the purpose of receiving
one reply or a narrowly bounded class of replies. In capability terms, it should
carry reply authority only, not general authority to invoke the client.

The book treats reply channels as a way to avoid leaking a client's full pid to
an untrusted peer. The general principle is that communication authority should
be as narrow and short-lived as the protocol step that requires it.

\paragraph{Service.}
A service is a durable, shared conversational resource intentionally offered by
a node for repeated use by multiple clients. In the Prolog Web, services may be
predicate-based or actor-based, but node-resident actor services are especially
important because they provide shared state or shared coordination over time.

In the current implementation, a published actor service is not just any actor
with a convenient name. It is something the node owner explicitly publishes
under a stable service name.

\paragraph{Server.}
A server, in the sense used by the book's generic-behavior chapters, is a
stateful actor that maintains private state and serves clients through a stable
request-response message protocol. The generic server behavior factors out the
machinery of message routing, state threading, and reply handling so that the
application-specific logic can be supplied as a callback predicate.

This is an Erlang-style behavior: it exposes a protocol, not a predicate-level
logical interface, and it does not bind the caller's variables directly.

\paragraph{Service Discovery.}
Service discovery is the process by which clients learn which services a node
intentionally makes available. In the Prolog Web, discovery should be treated as
an owner-curated publication mechanism, not as process introspection.

The important distinction is that discovery resolves published service names
without revealing everything that happens to be running. In the current book
example, this is realized declaratively through relations such as
\texttt{service/2} queryable via \texttt{rpc/2-3}.

\paragraph{Service Publication.}
Service publication is the act by which a node owner places a name into the
node's public service namespace. Publication is stronger than ordinary naming:
it makes a service part of the node's advertised public surface.

In the current regime, this is distinct from ordinary \texttt{register/2}.
Client registration is for client-managed actors, while published services are
managed through the owner-controlled service registry.

\paragraph{Session.}
A session is a bounded continuing interaction context in which state can persist
across multiple protocol steps. The book contrasts session-oriented interaction
with purely stateless request-response exchange.

In the current node, HTTP ISOTOPE sessions and WebSocket ACTOR sessions are both
session forms, but they have different timing and failure boundaries.

\paragraph{Shared Database.}
The shared database is the node-owned body of predicates and facts made
available as shared knowledge. It is loaded once at node startup into a shared
runtime module and then imported by actor modules or used by node-level query
execution.

Clients may query against shared knowledge and combine it with private code, but
the shared database is not the same thing as a client's private mutable
workspace.

\paragraph{Statechart Actor.}
A statechart actor is a Web Prolog actor whose control logic is given by a
statechart document. It behaves like an ordinary actor with a mailbox, but
incoming messages are interpreted as events that drive a structured state
machine rather than an ad hoc \texttt{receive/1-2} loop.

The distinctive semantic idea is run-to-completion: processing one external
event may trigger a finite cascade of internal work before the next external
message is consumed. This makes control explicit and inspectable while keeping
the process inside the ordinary actor model.

\paragraph{Supervisor.}
A supervisor is an actor whose main responsibility is to start, monitor, and
manage child actors, especially in the face of failure. In the book's Erlang-style
account, a supervisor embodies the ``let it crash'' discipline: workers are
allowed to fail, and the supervisor applies the configured recovery policy.

Supervisors matter most for durable node-resident services, where restart and
cleanup are part of the node owner's operational responsibility rather than
something each worker should reimplement ad hoc.

\paragraph{Toplevel Actor.}
A toplevel actor is an actor whose behavior is governed by an explicit
conversational protocol for submitting goals, receiving answers incrementally,
handling output, and managing abort or stop. It turns the ordinary Prolog
toplevel into a first-class actor process.

In implementation terms, toplevel actors underpin \texttt{rpc/2-3}, the
semi-stateful HTTP API, and much of the browser-facing shell machinery.

\paragraph{PTCP.}
PTCP is the Prolog Toplevel Communication Protocol. It is the explicit protocol
governing conversation between a client and a toplevel actor: how goals are
submitted, how answers are streamed incrementally, how output and prompts may
occur, and how abort, stop, and exit are handled.

Its importance in the book is that it turns informal REPL behavior into a
first-class behavioral contract. In implementation terms, PTCP underlies the
observable behavior of toplevel actors across the stateless, semi-stateful, and
stateful APIs.
